% Auto-generated by src/tables_figures/compute_portugal_validation.py
\begin{table}[htbp]
\centering
\caption{Directional validation against Portugal's 1996 reform (44$\to$40h).}\label{tab:portugal_validation}
\small
\begin{tabular}{lrrr}
\hline\hline
 & \textbf{Portugal 44$\to$40h} & \textbf{Model 44$\to$40h} & \textbf{Model 44$\to$36h} \\
 & \textit{Asai et al.\ (2025)} & \textit{This paper} & \textit{This paper, proposed} \\
 & firm-level DiD & aggregate & aggregate \\
\hline
$\Delta$Hours (\%) & $-10.9$ & $-3.71$ & $-8.71$ \\
$\Delta Y$ (\%) & $-3.2$ & $-2.00$ & $-8.02$ \\
$\Delta Y/h$ (\%) & $+4.4$ & $+1.77$ & $+0.75$ \\
$\Delta$Inf (pp) & \textit{n/a} & $+0.458$ & $+1.917$ \\
$A_{\text{req}}$ (\%) & \textit{n/a} & $+1.93$ & $+8.18$ \\
\hline\hline
\end{tabular}
\\[2pt]
\begin{minipage}{0.95\textwidth}
\justifying
{\footnotesize Signs match across all comparable outcomes. Magnitudes differ because the Portuguese estimates are firm-level, treated-versus-control partial-equilibrium estimates, while ours are economy-wide. The model's $e(h)$ captures only the fatigue channel, while the empirical productivity response also reflects worker selection, intensification, and sectoral capital deepening.}
\end{minipage}
\end{table}